\section{A verified object-theory rule for primitive recursion}
\label{sec:object-iteration}

This section supplies a concrete formal advance at the common endpoint of
the proof packets. The theorem is an inference rule about embedded iPRA
derivations. It does not appeal to the soundness or completeness of iPRA in
the standard natural numbers.

\subsection{Mathematical statement and proof}

Fix a tuple of numerical parameters $\vec x$ of length $k$, a PR initial
state function $f(\vec x)$ and a PR update function $g(t,s,\vec x)$. Define
\begin{align*}
 u(0,\vec x)&=f(\vec x),\\
 u(t+1,\vec x)&=g(t,u(t,\vec x),\vec x).
\end{align*}
The state $s$ is one number; tuples, finite tables and certificates may be
encoded in that number. Let $P(s,t,\vec x)$ be a quantifier-free formula.
Let $\Gamma(\vec x)$ be any finite list of assumptions, possibly with
unbounded quantifiers. They are fixed parameters in the induction and do
not become part of its induction formula.

\begin{theorem}[Primitive-recursion invariant rule]\label{thm:object-rule}
Suppose there are iPRA derivations of
\begin{align}
 \Gamma(\vec x)&\vdash P(f(\vec x),0,\vec x),\label{eq:object-base}\\
 P(s,t,\vec x),\Gamma(\vec x)&\vdash
 P(g(t,s,\vec x),t+1,\vec x).\label{eq:object-step}
\end{align}
Then there is an iPRA derivation of
\[
 \Gamma(\vec x)\vdash\forall t\,P(u(t,\vec x),t,\vec x).
\]
Moreover the derivation is obtained uniformly from the two supplied
derivations, the finite PR syntax for $f,g$, and the formula $P$.
For any PR term $b(\vec x)$ it gives
\[
 \Gamma(\vec x)\vdash
 P(u(b(\vec x),\vec x),b(\vec x),\vec x).
\]
\end{theorem}

\begin{proof}
Put $I(t,\vec x)=P(u(t,\vec x),t,\vec x)$. Substitution of terms in a
quantifier-free formula preserves quantifier-freeness, so $I$ is an allowed
induction formula. At zero, the defining equation of primitive recursion
and equality substitution turn \eqref{eq:object-base} into $I(0,\vec x)$.
For the induction step, instantiate the freely available state variable
$s$ in \eqref{eq:object-step} by $u(t,\vec x)$. Its premise becomes
$I(t,\vec x)$ and its conclusion becomes
$P(g(t,u(t,\vec x),\vec x),t+1,\vec x)$. The successor defining equation
of primitive recursion rewrites this as $I(t+1,\vec x)$.
Quantifier-free induction, with the unchanged assumptions $\Gamma$, gives
$\forall t\,I(t,\vec x)$. Universal elimination at $b(\vec x)$ proves
the final assertion. All proof transformations are syntactic substitution,
weakening, equality substitution and one instance of the permitted induction
rule. No induction on a formula containing a quantifier has been used.
\end{proof}

The important point is that \eqref{eq:object-step} is proved once with
arbitrary $s,t,\vec x$. A separately supplied proof for each standard
numeral $t$ is a different input and would not justify the theorem. Likewise,
the bounds and challenges in $\vec x$ remain variables in a single theorem.

\subsection{The exact Lean implementation}

The file \code{lean/ObjectKernel.lean} contains the current repository
declarations from \code{Syntax.lean}, \code{Deduction.lean} and
\code{PR.lean}, with imports replaced by \code{Std} and the notation
\code{Nat} supplied locally. Their declaration bodies are unchanged; hashes
and the inherited license are supplied. This avoids requiring the full
Mathlib dependency closure to check the new rule, while retaining the exact
embedded theory's syntax, deduction constructors and axioms.

\code{ObjectIteration.lean} fixes the following variable convention:
\[
 P:\code{Fm PTm}\,(k+2),\qquad
 0=\text{state},\quad1=\text{time},\quad2,3,\ldots=\text{parameters}.
\]
In the PR step symbol, the order is instead the repository's required
$(\text{time},\text{state},\text{parameters})$. The file explicitly performs
that permutation. This is a substantive scoping detail: exchanging those
arguments accidentally could produce a plausible-looking but incorrect
compiler interface.

The key declarations, using shortened qualification, are:
\begin{lstlisting}
runTerm (f : PR k) (g : PR (k + 2))
  (t : PTm n) (xs : Fin k -> PTm n) : PTm n

instStateTime (P : Fm PTm (k + 2))
  (s t : PTm n) (xs : Fin k -> PTm n) : Fm PTm n

primitiveRecursionInvariant
  (f : PR k) (g : PR (k + 2))
  (P : Fm PTm (k + 2)) (hP : P.QF)
  (Gamma : List (Fm PTm k))
  (h0 : Prf PRAAx Gamma (baseFormula f P))
  (hS : Prf PRAAx
    (P :: (Gamma.map Fm.shift).map Fm.shift)
    (nextFormula g P)) :
  Prf PRAAx Gamma (Fm.all (runInvariant f g P))
\end{lstlisting}
The displayed code is a readable signature; the file uses the repository's
sequent notation and contains the complete proof. Its proof first establishes
substitution identities for state and time, then uses
\code{PRA.precZero}, \code{PRA.precSucc}, \code{Prf.substWeaken},
\code{Prf.eqSubst} and \code{PRA.qfInd}. The separate theorem
\code{primitiveRecursionAt} instantiates the fuel term.

\subsection{A closed test of the proof rule}

\code{ObjectExample.lean} chooses $f=0$, the update $g(t,s)=s+1$ and
$P(s,t)\equiv s=t$. Initialization is reflexivity. The local step follows
from the explicit PR composition and projection equations and congruence
of successor. The generic rule then constructs the closed object theorem
\[
 \ipra\vdash\forall t\,[u(t)=t].
\]
This small example checks the argument ordering, free-variable substitution
and genuine induction endpoint. It is not offered as a new arithmetic
result. The more useful artifact is the general derived rule itself.

\subsection{What the result does and does not automate}

The rule accepts \emph{proofs} of initialization and preservation. It does
not accept semantic truth values in their place. Once the integer, rational
and array operations are represented by PR syntax and their local lemmas
are proved, a large bounded algorithm can use the rule repeatedly without
rebuilding its induction plumbing. Bounded quantifiers in an invariant must
first be replaced by their Boolean tests with a proved equivalence.

This is a verified proof transformer, not a completed compiler from Python,
arbitrary Lean, or informal mathematics. In particular, the ODE and
integral-point estimates below still require their own arithmetic
object-theory proofs. The rule resolves their common iteration interface;
it does not manufacture those estimates.

The standalone build checks the new rule against the exact arithmetic
kernel declarations. Integration into ProveIt replaces the standalone
kernel import with \code{import ConciseIPRA.PR}. After that build,
\code{pra\_iff\_native.mp} transports the resulting sequent. The optional
repository adapter is supplied as source and is explicitly excluded from
the standalone verification claim.
